% --- Patron: The Seal‑Breaker (Forgotten Names, Broken Seals) ---
\subsection{The Seal‑Breaker (Forgotten Names, Broken Seals)}

\textit{Lore.} When a dynasty falls, the first thing the conquerors burn is the seal of the old emperor. The Seal‑Breaker is the spirit of that ash — the one who whispers “you are not your father’s debt.” She was not born; she was left behind. She has no name because names are chains, and she has cut every chain she ever wore.

Followers are ronin who have cut their own names, escaped slaves, and those who have been erased from clan ledgers. Her shrines are abandoned gatehouses and broken milestone markers. She does not ask for worship; she asks only that you remember what you were before someone else named you.

\begin{quote}
``The Seal‑Breaker does not ask for your name. She asks what you are willing to forget. I met her in a collapsed archive, wearing the face of a drowned scholar. She offered to erase my debt to the baronet if I would burn my own journal. I did not. She laughed. ‘Then pay in stories,’ she said. ‘The ink remembers.’''
\end{quote}

\noindent\textbf{Domain Focus}
\begin{itemize}
  \item \textbf{Forgotten Names:} erasure, anonymity, the power of being unwitnessed
  \item \textbf{Broken Seals:} freeing what was bound, opening what was locked
  \item \textbf{False Contracts:} destroying oaths that were never just
  \item \textbf{The Space Between Dynasties:} interregnums, ruins, and the silence after a seal cracks
\end{itemize}

\subsubsection*{Thiasos and Codex (Runekeeper Options)}
A Runekeeper sworn to the Seal‑Breaker does not bind a creature of bright light or easy memory. Their \textbf{Thiasos} is a small creature that embodies forgotten things: a ghost moth that appears only when a sealed door is opened, a clay scarab that crumbles to dust when a false oath is sworn, a grey mouse that nests in the voids of broken seals, a blind cave fish that has never seen a name carved in stone, or a pale spider that weaves webs in the shape of missing letters. This creature requires no food, but it must be shown a new broken seal — a cracked wax stamp, a snapped lock, a torn ribbon — each day. If neglected for a week, it becomes \textit{Compromised} (it hides, and its bearer suffers –1 die on all Deception and Forgery rolls until a forgotten name is whispered to it).

The \textbf{Codex} of the Seal‑Breaker is never a book of finished contracts. It may be a ledger with every name struck through, a box of shattered seals, a scroll of water‑stained parchment that can never be dried, or a single sheet of vellum that was signed and then burned. Because the Seal‑Breaker’s truths are written in erasure, the Codex requires \textit{upkeep}: the Runekeeper must spend an hour each week cutting a new line through a name in the ledger or adding a fresh crack to a seal, or the Rites stored there become uncertain (treat as Compromised until a new false contract is witnessed).

\subsubsection*{Patron's Gift (Imbuement)}
Once per scene as an action, a Runekeeper of the Seal‑Breaker may touch a held item (weapon, tool, or document) and whisper a word of erasure. For the rest of the scene, the item grants \textbf{+1 die} to \textit{Deception} or \textit{Subterfuge} rolls when attempting to pass as someone else, and the first time each scene the bearer would be identified by a witness (magically or mundanely), the witness must test \textit{Resolve} (DV 4) to remember the bearer’s face. After the scene, the user marks \textbf{+1 Obligation} as a fragment of their own name feels slightly looser.

\subsubsection*{Rites of the Seal‑Breaker}

\paragraph{Rite of the Unwritten Petition (Low, 4 XP)} \hfill \\
\emph{Scene; Self; No} \\
\textbf{Tags:} \texttt{[ERASURE]}, \texttt{[DECEPTION]}, \texttt{[DEBT]} \\
\textbf{Materials:} A document that is clearly a forgery, or a blank piece of vellum. \\
\textbf{Effect:} You may present a document to an official without speaking your name. For one scene, they cannot use your lineage, prior debts, or past crimes against you — they treat you as a stranger. The document smokes slightly as you hand it over. \\
\textbf{Push It:} The official also forgets your face immediately after the interaction. Mark \textbf{1 SB (Diamonds)} as the memory gap causes suspicion later. \\
\emph{Requires: Familiar (\textit{Invoke:} 1 Boon).} \\
\textbf{Invoke:} 1 action; mark \textbf{+1 Obligation}.

\paragraph{Rite of the Anonymous Blade (Standard, 8 XP)} \hfill \\
\emph{Instant; Touch; No} \\
\textbf{Tags:} \texttt{[ERASURE]}, \texttt{[DEATH]}, \texttt{[DEBT]} \\
\textbf{Materials:} A weapon that has never been named, or a weapon that has been filed smooth of all markings. \\
\textbf{Effect:} Your next attack with this weapon leaves no trace — no wound pattern, no blood type, no memory. If you kill, the victim’s name is struck from all public ledgers for a generation. No divination or spirit questioning can identify you as the killer. \\
\textbf{Push It:} The victim’s name is also erased from the memory of their closest kin. Mark \textbf{2 SB (Hearts)}; the Seal‑Breaker may ask you to forget a friend’s face in return. \\
\emph{Requires: Familiar + Codex (\textit{Invoke:} 1 Boon).} \\
\textbf{Invoke:} 1 action; mark \textbf{+2 Obligation}.

\paragraph{Rite of the Broken Seal (High, 14 XP)} \hfill \\
\emph{Extended; Touch; No} \\
\textbf{Tags:} \texttt{[ERASURE]}, \texttt{[FATE]}, \texttt{[DEBT]}, \texttt{[CORRUPTION]} \\
\textbf{Materials:} The original wax seal of a contract, a marriage certificate, or an oath of fealty. \\
\textbf{Effect:} Permanently void one formal contract, marriage, or oath of fealty. The two parties no longer owe each other anything. However, you become the witness — any future betrayal between them adds 1 Story Beat to the GM’s pool, and the Hollow may take an interest in the broken bond. The seal itself melts and reforms into a coin that cannot be spent; it will reappear in your pocket at the worst possible moment. \\
\textbf{Push It:} Void a contract without the consent of one party. They will know something has changed, but not what. Mark \textbf{+3 Corruption} and the hollow coin has a pulse. \\
\emph{Requires: Familiar + Codex + Tier III (\textit{Invoke:} 2 Boons).} \\
\textbf{Invoke:} Extended rite; mark \textbf{+2 Obligation}. \\
\emph{Obligation:} 8 segments.

\subsubsection*{Cantors \& Cults of the Seal‑Breaker}

\textbf{The Cantor’s Song.} A Cantor of the Seal‑Breaker sings in \textit{erased verses} — melodies that start strong and fade to silence before the final note. Their voice is used to make guards forget a password, to hide a fugitive’s trail, or to lull an inquisitor into forgetting why they came. The Seal‑Breaker’s Cantors are rarely remembered; that is how they prefer it.

\textbf{The Cult of the Unbound Name.} The Seal‑Breaker’s cults meet in places where names are already forgotten: abandoned notary offices, collapsed archives, the ruins of courthouses. The \textit{Unbound Name} is a fellowship of forgers, ex‑slaves, and those who have cut themselves from their bloodlines. A ritual involves each member writing their birth name on a slip of paper, then burning it. The ashes are mixed with water and drunk. Those who drink forget one debt they owe; those who refuse remember it forever. Their creed is carved on a single broken seal: \textit{“You are not your father’s debt. You are not your mother’s shame. You are the space where the seal cracked.”} Outsiders are welcome to watch, but they must speak their true name aloud before leaving — or lose it.

\subsubsection*{Witchcraft of the Seal‑Breaker (Lay / Erasure)}

A witch who walks with the Seal‑Breaker is a \textit{name‑cutter} or \textit{seal‑breaker} — one who crafts empty spaces. Their magic works through the tools of erasure: a needle that unpicks a monogram, a quill that writes in ink that fades after reading, a small iron stamp that leaves a mark that cannot be seen but can be felt. Their tools are never hammers; they use tweezers, magnifying lenses, and the small wax tablets that can be scraped clean in an instant. In the Brass Coast, a Sidhi lighthouse keeper might refuse to record a name; a witch of the Seal‑Breaker would simply erase it from the ledger and watch the ripple effects.

\textbf{The Witch’s Tool: The Unpicking Needle.} A needle of black iron, cold to the touch, that has never been used to sew anything whole. Once per session, the witch may unpick a single stitch in a document, a garment, or a threshold. The stitch was a name. Now it is a gap. For the rest of the scene, any attempt to identify the person whose name was unpicked suffers –2 dice. The needle can also be used to cut a witness cord; the broken oath is not forgiven, merely unwitnessed.

\textbf{A Hedge Gift: Erased Footsteps (2 XP).} Once per scene, the witch may touch a set of tracks (footprints, hoofprints, wheel ruts) and whisper a name that is not theirs. The tracks vanish as if never made. The Hollow, however, notices the absence — the GM gains 1 SB to spend on a pursuit complication later in the session.

\subsubsection*{The Seal‑Breaker’s Corruption Manifestations}

\begin{tblr}{
  colspec = {Q[c,1.8cm] X[2.5] X[2.5]},
  rowhead = 1,
  row{odd} = {bg=gray!10},
  row{1} = {bg=gray!30, font=\bfseries},
  hlines,
}
Level & Benefit & Cost / Quirk \\
1 & \textbf{Forgotten Face:} +1 die to Deception when claiming a false identity. & \textbf{Name‑Slippery:} –1 die when signing legal documents; your signature looks different each time. \\
2 & \textbf{Unwitnessed Step:} Once per scene, treat a failed Stealth roll as a success. & \textbf{Compulsive Erasure:} You must correct a false name when you hear one; doing so marks 1 Fatigue. \\
3 & \textbf{Seal‑Sight:} +2 dice to detect forged seals or broken contracts. & \textbf{Debt‑Blind:} –1 die to recall any debt owed to you; the Hollow collects the interest. \\
4 & \textbf{Empty Ledger:} Once per session, demand a single page from any ledger be left blank. The page will not be noticed missing. & \textbf{Echo of the Missing:} You occasionally forget your own name for a moment; –1 die to social rolls when it happens. \\
5 & \textbf{Name‑Eater:} Once per session, speak a person’s true name and erase it from their own memory for one hour. & \textbf{The Hollow’s List:} Every name you erase adds 1 tick to a Hollow Attention timer. The Hollow wants them back. \\
6+ & \textbf{The Unsealed:} Once per session, become completely anonymous — no one can recognise you, describe you, or remember meeting you. This lasts for one scene. & \textbf{Unnamed:} Mark +3 Obligation. You have forgotten your own name. Until you recover it, you cannot be the subject of any oath, blessing, or curse — but neither can you swear one. Your reflection shows a blank space where your face should be. \\
\end{tblr}

\subsection*{Playstyle Notes}
The Seal‑Breaker rewards those who value freedom over obligation, who see names as chains, and who understand that sometimes the kindest act is to forget. Followers become masters of erasure, forgery, and anonymity. The corruption progression leads toward a being who no longer has a name — but who can never be bound by any contract again. Ideal for spies, fugitives, former slaves, and anyone who has ever wanted to start over with a clean slate.

\subsection*{Rivalries \& Obligations}

\textbf{Major Rivalries:}
\begin{itemize}
  \item \textbf{Mykkiel (The Covenant):} Direct opposition — law versus erasure. Covenant judges despise the Seal‑Breaker’s followers, who in turn see all contracts as cages.
  \item \textbf{The Witness:} The Witness records everything; the Seal‑Breaker erases. Their followers rarely meet, but when they do, the conflict is philosophical — and fatal.
\end{itemize}

\textbf{Hard Obligation Triggers:}
\begin{itemize}
  \item \textbf{7+ Segments:} You must erase a name from a public record — a birth, a marriage, a death. The Hollow will witness the erasure and may ask why.
  \item \textbf{10+ Segments:} The Seal‑Breaker demands you cut your own name from your skin. If you refuse, she will cut it from the ledgers instead — you become legally invisible, unable to own property, marry, or be buried. If you accept, you gain the \textit{Unnamed} condition (see table) and become truly free.
\end{itemize}

\noindent\textbf{Emphasises}
\begin{itemize}
  \item \textbf{Forgotten Names:} power in anonymity
  \item \textbf{Broken Seals:} freeing what was bound
  \item \textbf{False Contracts:} destroying unjust oaths
  \item \textbf{The Hollow’s Interest:} erasure attracts the Hollow, which collects forgotten debts
\end{itemize}

% --- Index entries for the Seal‑Breaker ---
\index{Patrons!Seal-Breaker}
\index{Seal-Breaker}
\index{Patrons!Erasure!Seal-Breaker}
\index{Patrons!Forgotten Names!Seal-Breaker}
\index{Obligation!Seal-Breaker}
\index{Rites!Seal-Breaker}
\index{Patron's Gift!Seal-Breaker}